\documentclass[11pt]{article}
\usepackage[margin=25mm]{geometry}
\usepackage{fontspec}
\setmainfont{Latin Modern Roman}
\usepackage{amsmath,amssymb,mathtools}
\usepackage{unicode-math}
\setmathfont{Latin Modern Math}
\usepackage{xurl,hyperref}
\hypersetup{colorlinks=true,urlcolor=blue}
\setcounter{secnumdepth}{0}
\setlength{\parindent}{0pt}
\setlength{\parskip}{5pt}
\setlength{\emergencystretch}{3em}
\newcommand{\sref}[2]{\hyperlink{source:#1:#2}{[S#2]}}
\newcommand{\cataloguescope}[1]{\paragraph{Recorded target.}#1}
\newcommand{\R}{\mathbb R}
\newcommand{\C}{\mathbb C}
\newcommand{\N}{\mathbb N}
\newcommand{\Z}{\mathbb Z}
\newcommand{\Q}{\mathbb Q}
\newcommand{\F}{\mathbb F}
\newcommand{\eps}{\varepsilon}
\begin{document}
% UnsolvedMath ID 30000578; source catalogue code OWR-1325-003
\setcounter{section}{30000577}
\section{Topographic Rossby Waves and Inertial-Wave Coexistence}\label{30000578}
{\small UnsolvedMath ID 30000578 · Mathematical Physics}
\subsection{Definitions and mathematical statement}

\paragraph{Shared notation.}$\mathbb N=\{1,2,\ldots\}$; the standard number fields have their usual meanings. All additional conventions are stated locally.


\paragraph{Definitions and conventions.} A rotating fluid supports inertial waves due to Coriolis restoring forces. A sloping boundary or varying depth gives an effective topographic potential-vorticity gradient and supports topographic Rossby waves. The report compares these waves in a rectangular rotating basin with a sloping boundary and examines an asymptotic series whose leading term is a topographic Rossby wave.

\paragraph{Source scope and qualifications.} The basin geometry, boundary slope, and asymptotic expansion of the contribution are part of the context. The source does not ask whether every Rossby wave in every geometry diverges; the two alternatives concern the particular paradox described there.

\paragraph{Catalogue mathematical target.}
\paragraph{Catalogue wording retained for traceability.}
How is the apparent paradox between topographic Rossby waves and inertial waves resolved: does the series whose zeroth-order term is the topographic Rossby wave diverge locally, or does it coexist with the inertial wave?

\paragraph{Saved literature context.} The supplied catalogue's assessment dated 2026-08-21 records: The report presents this as an open problem, conjecture, or unresolved research direction; targeted literature checking found no authoritative resolution of the extracted claim as of 2026-08-21. The assessment is retained as dated provenance; it is not a new comprehensive literature or open-status certification.
\subsection{Short English statement}
Explain whether the topographic-Rossby expansion breaks down locally or describes waves coexisting with the exact inertial waves.
\subsection{Research attempt}
\paragraph{Model, date and investigated scope.}
OpenAI GPT-6 Astra (configured model \texttt{gpt-6-astra}), ultra reasoning; generated 11 October 2026. Explicit two-mode calculation shows coexistence of exact branches can accompany failure of a uniform slow-mode series.

\paragraph{Sources and literature check.}
The complete supplied definition, including source qualifications and the dated literature assessment, was reviewed. The calculation below is an elementary restricted attempt motivated by [S2]; it does not claim novelty or a fresh literature-status determination.

\paragraph{Context and investigated scope.} The complete supplied definition, including source qualifications and the dated literature assessment, was reviewed. The calculation below is an elementary restricted attempt motivated by [S2]; it does not claim novelty or a fresh literature-status determination.

\paragraph{A two-branch singular-perturbation test.} Let $\omega>0$ represent an inertial frequency and let $\epsilon$ measure coupling to a slow branch. The real symmetric matrix
\[
 L_{\epsilon,\omega}=\begin{pmatrix}0&\epsilon\\\epsilon&\omega\end{pmatrix}
\]
has eigenvalues $\lambda_\pm=(\omega\pm\sqrt{\omega^2+4\epsilon^2})/2$. For $|\epsilon|<\omega/2$, the slow eigenvalue expands as $\lambda_-=-\epsilon^2/\omega+\epsilon^4/\omega^3+\cdots$. The Taylor series has radius $\omega/2$, since the square root has branch points at $\epsilon=\pm i\omega/2$.

\paragraph{Coexistence does not guarantee a uniform series.} For each real $\epsilon$ both exact eigenvalues exist and are distinct unless $\epsilon=\omega=0$. If a sequence of inertial modes has frequencies $\omega_n\downarrow0$, the radii of convergence shrink to zero. Thus no fixed nonzero coupling admits this expansion uniformly in $n$, even though every exact two-mode problem has two coexisting branches. At $\omega=0$ the ordered eigenvalues are $\pm|\epsilon|$, and the expansion from fixed positive $\omega$ plainly fails.

\paragraph{Limit of the analogy.} This does not identify $L_{\epsilon,\omega}$ with the rectangular sloping-basin operator or prove local divergence of its Rossby series. It demonstrates that the source's two alternatives need not be mutually exclusive at the level of perturbation reasoning: exact modes can coexist while a mode-uniform expansion fails. The first missing step is a spectral reduction of the actual boundary-value problem, including its small divisors and spatial convergence norm. Without that reduction the observed basin paradox remains unresolved.

\paragraph{Verification and adversarial audit.}
Checked eigenvalues, complex branch points, shrinking radius and zero-frequency degeneracy; no claim of an actual basin counterexample. Model self-audit only; no independent human proof checking. Prompts A and B were applied to this canonical record and its qualified source scope. Checked eigenvalues, complex branch points, shrinking radius and zero-frequency degeneracy; no claim of an actual basin counterexample. Retained unresolved status and zero original-target progress; The actual sloping-basin operator has not been reduced with quantitative small-divisor and spatial convergence control.

Coordinator checked the two eigenvalues, expansion coefficients and shrinking complex convergence radius as the gap tends to zero. This additional check was model-assisted, not independent human verification.

\paragraph{Tools, process and contributions.}
Python 3 was used for file preparation and structural checks; no numerical experiment or theorem prover was used. Contribution prompt A was followed by prompt B. GPT-6 Astra with ultra reasoning drafted the research and performed its stated self-audit; any additional coordinator contribution is identified in the verification and audit. A GPT-6-based Codex coordinator assembled the manuscript and checked repository integration; its exact variant and reasoning configuration were not exposed. No human mathematical contribution, independent human proof check or formal proof certification is supplied. Complete derivations and the consulted source references are included here; private reasoning transcripts are not published.

\paragraph{Final status and first remaining gap.}
\textbf{UNRESOLVED}. The actual sloping-basin operator has not been reduced with quantitative small-divisor and spatial convergence control.

\section{Completion Estimate}
\noindent\textbf{COMPLETION: 0\%}
The calculation is an auxiliary or restricted attempt and establishes no newly verified portion of the full source target. The actual sloping-basin operator has not been reduced with quantitative small-divisor and spatial convergence control.

\subsection{Sources}
\begin{itemize}
\item[\hypertarget{source:30000578:1}{S1}] ULAM.AI and the UnsolvedMath Contributors, UnsolvedMath, record 30000578 (OWR-1325-003). Catalogue curation under CC BY 4.0; definitions and English exposition added locally.
\url{https://huggingface.co/datasets/ulamai/UnsolvedMath}.
\item[\hypertarget{source:30000578:2}{S2}] Mathematical Theory and Modelling in Atmosphere-Ocean Science (2006), Oberwolfach Reports, DOI 10.4171/owr/2006/39, pp. 2307--2398; official report source and mathematical context.
\url{https://doi.org/10.4171/owr/2006/39}.
\item[\hypertarget{source:30000578:3}{S3}] Official EMS Press PDF of the same Oberwolfach report; PDF page 75.
\url{https://ems.press/content/serial-article-files/46068}.
\end{itemize}
\label{end:30000578}
\end{document}
